% Generated draft: requires mathematical and source review.
\documentclass[11pt]{article}
\usepackage[margin=25mm]{geometry}
\usepackage{fontspec}
\setmainfont{Latin Modern Roman}
\usepackage{amsmath,amssymb,mathtools,unicode-math}
\setmathfont{Latin Modern Math}
\usepackage{enumitem,needspace,xurl,hyperref}
\setlength{\parindent}{0pt}
\setlength{\parskip}{5pt}
\setlength{\emergencystretch}{3em}
\newcommand{\N}{\mathbb{N}}
\newcommand{\Z}{\mathbb{Z}}
\newcommand{\Q}{\mathbb{Q}}
\newcommand{\R}{\mathbb{R}}
\newcommand{\C}{\mathbb{C}}
\newcommand{\F}{\mathbb{F}}
\newcommand{\A}{\mathbb{A}}
\newcommand{\PP}{\mathbb{P}}
\newcommand{\E}{\mathbb{E}}
\newcommand{\eps}{\varepsilon}
\newcommand{\dd}{\,\mathrm{d}}
\begin{document}
\section*{UnsolvedMath ID: 30004518}
\textbf{Source problem code:} catalogue entry
\subsection{Definitions and mathematical statement}
Let \(n\geq 3\) be an integer and let \(A=(\alpha_0,\ldots,\alpha_{n-1})\) be an ordered angle sequence. A polygonal realization is understood to be closed and to carry the prescribed angles in their cyclic order. The supplied record does not define the angle convention, the allowed range of the \(\alpha_i\), or the admissible degeneracies; these remain unresolved qualifications.

Write \(\mathcal R(A)\) for the assertion that \(A\) is realizable by a polygon in \(\R^3\), and \(\mathcal S(A)\) for the assertion that such a realization exists without self-intersections. For a nondegenerate polygon, absence of self-intersections means that its closed polygonal parametrization is injective apart from the identification of its initial and final parameter values. Whether the source permits degeneracies or uses additional conventions is unspecified.

The proposed assertion, with the source's intended conventions still to be supplied, is that for every integer \(n\geq 3\) and every ordered angle sequence \(A\) satisfying \(\mathcal R(A)\),
\[
\mathcal S(A)
\quad\Longleftrightarrow\quad
\left(n\text{ is even}\right)
\ \lor\ 
\left(\sum_{i=0}^{n-1}(\pi-\alpha_i)\neq\pi\right).
\]
The displayed criterion is retained exactly as supplied. A fully specified mathematical statement requires clarification of the angle and degeneracy conventions.
\subsection{Short English statement}
Given an angle sequence already realizable by a closed polygon in \(\R^3\), the question asks whether a realization without self-intersections exists exactly when the number of angles is even or the displayed sum differs from \(\pi\). The source does not specify the angle convention or allowed degeneracies. Its open-status label and literature assessment are supplied curation requiring independent review.
\subsection{Sources}
\begin{itemize}
\item Supplied local UnsolvedMath registry record, permanent ID \(30004518\), title \emph{Simple Spatial Polygon Realizations from Angle Sequences}, source problem code \texttt{catalogue entry}, in the Oberwolfach Reports Open Problems collection. No individual proposer is identified; the recorded proposed year is 2020.
\item Supplied source citation: \emph{Discrete Geometry (2021)}, Oberwolfach Reports, DOI 10.4171/owr/2020/30, pp.\ 1469--1528. \url{https://doi.org/10.4171/owr/2020/30}. The record describes Oberwolfach Reports as documenting mathematical workshops and their open-problem discussions.
\item Supplied literature reference: Efrat--Rote--Rutter--Shabbeer, \emph{Polygons with Prescribed Angles}. \url{https://jgaa.info/index.php/jgaa/article/view/paper599}.
\item Supplied literature reference: \emph{Oberwolfach Reports 2020/30}. \url{https://ems.press/content/serial-article-files/46867}.
\item Literature assessment dated exactly \texttt{2026-08-22}: the supplied curation labels the problem open and describes the equivalence as conjectural, with only its necessary direction and realizability algorithms proved. This is supplied dated curation derived from the OWR record, not an independently verified assessment; time-sensitive status and cited claims require independent verification.
\item The registry labels the statement \texttt{corrected\_verified}, explaining that only formula delimiters, notation, grammar, and typography were repaired and that mathematical content and scope were preserved. This supplied editorial label does not establish mathematical correctness.
\end{itemize}
\end{document}
